\section{Introduction and problem statement}\label{sec:intro}

\paragraph{Motivation.}
Questions about long videos are rarely answered by a single moment. ``What
colour was the board before it was painted?'' requires the system to have
noticed the board early on, to have registered the change later, and to
recognise that both observations concern the same object. ``Which tool used at
the beginning reappears at the end?'' requires linking entities across tens of
minutes. Current vision-language models (VLMs) process a bounded number of
visual tokens; a 60-minute video sampled at even one frame per second exceeds
that bound by orders of magnitude, and uniform subsampling discards exactly the
fine, transient details such questions target.

The alternative studied here is to observe the video incrementally, in short
chunks, and to maintain a \emph{memory}: a compact, queryable representation
written as the video unfolds and consulted when a question is asked. Because
the memory must itself fit into the answering model's context, it operates
under a \emph{budget}. The system must therefore decide what to store, how to
update stored knowledge when the world changes, which mentions refer to the
same entity, what to compress and what to forget.

\paragraph{Problem statement.}
Given one long video, a fixed perception pipeline that converts each chunk into
an observation, and a fixed memory budget, we ask which \emph{organisation} of
memory best supports answering questions about past events, changes of
objects, and relations between distant parts of the video, and whether the
answers it supports can be traced back to the correct source clips.

The object of study is the memory organisation, not the perception model. This
has a strict methodological consequence that shapes the whole design: every
memory variant must consume \emph{identical} per-chunk observations and operate
under the \emph{same} budget, measured in the same unit, and must be queried by
the same answering model with the same retrieval allowance. Differences in
accuracy can then be attributed to memory organisation alone. The pipeline is
therefore built in stages whose outputs are hashed and frozen before the next
stage starts, so any later deviation from them is detectable.

\paragraph{Contributions.}
\begin{enumerate}[nosep]
  \item A formal statement of budgeted incremental video memory and of the
        evaluation quantities (\cref{sec:formulation}).
  \item Four research questions derived from the project brief, each with
        operational measures (\cref{sec:rq}).
  \item A system design whose two invariants, memory-agnostic perception and
        provenance-preserving, event-sourced memory, make the planned
        comparisons valid (\cref{sec:overview}).
  \item A documented pipeline in five stages: environment check,
        video probing and reproducible preprocessing
        (\cref{sec:stage1}); a frozen evaluation set (\cref{sec:stage2});
        audited VLM perception (\cref{sec:stage3}); a structured memory with
        entity resolution, debounced state changes, compression and
        priority-based forgetting, with three baselines (\cref{sec:stage4});
        and an evaluation protocol with failure attribution, audits and an
        interactive demo (\cref{sec:stage5}).
\end{enumerate}
